% GENERATED FILE. Edit canonical lessons, then run npm run generate:books.
% canonical-source-hash: fnv1a64:456d9817bd2557b3
% canonical-lessons: TE-C306-ukku, TE-C306-simentu, TE-C306-tolu, TE-C306-patti, TE-C306-unni

\chapter{Steel, Cement, Leather, Cotton, Wool}
\label{ch:te-steel-cement-leather-cotton-wool}
\hlchaptermodality{306}

\begin{chapteropening}
\textbf{By the end of this chapter:} \emph{I can say steel, cement, leather, cotton and wool.}

Complete the last lesson of chapter 306: I can say steel, cement, leather, cotton and wool.
\end{chapteropening}

\section[\texorpdfstring{\te{ఉక్కు} (ukku)}{ఉక్కు (ukku)}]{\te{ఉక్కు} (ukku) — steel}

\label{lesson:TE-C306-ukku}



\begin{quote}
Before the new one: say the Telugu for a battery, then the Telugu for a video.
\end{quote}

\subsection*{You'll want to know: \te{ఉక్కు}}
\textbf{\te{ఉక్కు}} — \emph{ukku} — \textquotedblleft{}steel\textquotedblright{}.

\begin{grammarlens}[title={What you've built}]
One more word for this chapter.
\end{grammarlens}

\subsection*{Your turn}
\begin{itemize}
\raggedright
  \item \emph{Say it:} \emph{ukku}
  \item \emph{Say it:} \emph{ukku}, once more
  \item \emph{Say it:} say \emph{ukku}
  \item \emph{Recall:} say the Telugu for a battery, then the Telugu for a video, then say \emph{ukku} again
\end{itemize}

\begin{tcolorbox}[breakable,colback=teal!4,colframe=teal!35!black,title={Before you move on}]
What does \te{ఉక్కు} mean? (\textquotedblleft{}Steel\textquotedblright{}.) Say it once more.
\end{tcolorbox}
\section[\texorpdfstring{\te{సిమెంటు} (simeṇṭu)}{సిమెంటు (simeṇṭu)}]{\te{సిమెంటు} (simeṇṭu) — cement}

\label{lesson:TE-C306-simentu}



\begin{quote}
Before the new one: say the Telugu for a video, then the Telugu for steel.
\end{quote}

\subsection*{You'll want to know: \te{సిమెంటు}}
\textbf{\te{సిమెంటు}} — \emph{simeṇṭu} — \textquotedblleft{}cement\textquotedblright{}.

\begin{grammarlens}[title={What you've built}]
One more word for this chapter.
\end{grammarlens}

\subsection*{Your turn}
\begin{itemize}
\raggedright
  \item \emph{Say it:} \emph{simeṇṭu}
  \item \emph{Say it:} \emph{simeṇṭu}, once more
  \item \emph{Say it:} say \emph{simeṇṭu}
  \item \emph{Recall:} say the Telugu for a video, then the Telugu for steel, then say \emph{simeṇṭu} again
\end{itemize}

\begin{tcolorbox}[breakable,colback=teal!4,colframe=teal!35!black,title={Before you move on}]
What does \te{సిమెంటు} mean? (\textquotedblleft{}Cement\textquotedblright{}.) Say it once more.
\end{tcolorbox}
\section[\texorpdfstring{\te{తోలు} (tōlu)}{తోలు (tōlu)}]{\te{తోలు} (tōlu) — leather}

\label{lesson:TE-C306-tolu}



\begin{quote}
Before the new one: say the Telugu for steel, then the Telugu for cement.
\end{quote}

\subsection*{You'll want to know: \te{తోలు}}
\textbf{\te{తోలు}} — \emph{tōlu} — \textquotedblleft{}leather\textquotedblright{}.

\begin{grammarlens}[title={What you've built}]
One more word for this chapter.
\end{grammarlens}

\subsection*{Your turn}
\begin{itemize}
\raggedright
  \item \emph{Say it:} \emph{tōlu}
  \item \emph{Say it:} \emph{tōlu}, once more
  \item \emph{Say it:} say \emph{tōlu}
  \item \emph{Recall:} say the Telugu for steel, then the Telugu for cement, then say \emph{tōlu} again
\end{itemize}

\begin{tcolorbox}[breakable,colback=teal!4,colframe=teal!35!black,title={Before you move on}]
What does \te{తోలు} mean? (\textquotedblleft{}Leather\textquotedblright{}.) Say it once more.
\end{tcolorbox}
\section[\texorpdfstring{\te{పత్తి} (patti)}{పత్తి (patti)}]{\te{పత్తి} (patti) — cotton}

\label{lesson:TE-C306-patti}



\begin{quote}
Before the new one: say the Telugu for cement, then the Telugu for leather.
\end{quote}

\subsection*{You'll want to know: \te{పత్తి}}
\textbf{\te{పత్తి}} — \emph{patti} — \textquotedblleft{}cotton\textquotedblright{}.

\begin{grammarlens}[title={What you've built}]
One more word for this chapter.
\end{grammarlens}

\subsection*{Your turn}
\begin{itemize}
\raggedright
  \item \emph{Say it:} \emph{patti}
  \item \emph{Say it:} \emph{patti}, once more
  \item \emph{Say it:} say \emph{patti}
  \item \emph{Recall:} say the Telugu for cement, then the Telugu for leather, then say \emph{patti} again
\end{itemize}

\begin{tcolorbox}[breakable,colback=teal!4,colframe=teal!35!black,title={Before you move on}]
What does \te{పత్తి} mean? (\textquotedblleft{}Cotton\textquotedblright{}.) Say it once more.
\end{tcolorbox}
\section[\texorpdfstring{\te{ఉన్ని} (unni)}{ఉన్ని (unni)}]{\te{ఉన్ని} (unni) — wool}

\label{lesson:TE-C306-unni}



\begin{quote}
Before the new one: say the Telugu for leather, then the Telugu for cotton.
\end{quote}

\subsection*{You'll want to know: \te{ఉన్ని}}
\textbf{\te{ఉన్ని}} — \emph{unni} — \textquotedblleft{}wool\textquotedblright{}.

\begin{grammarlens}[title={What you've built}]
One more word for this chapter.
\end{grammarlens}

\subsection*{Your turn}
\begin{itemize}
\raggedright
  \item \emph{Say it:} \emph{unni}
  \item \emph{Say it:} \emph{unni}, once more
  \item \emph{Say it:} say \emph{unni}
  \item \emph{Recall:} say the Telugu for leather, then the Telugu for cotton, then say \emph{unni} again
\end{itemize}

\begin{tcolorbox}[breakable,colback=teal!4,colframe=teal!35!black,title={Before you move on}]
What does \te{ఉన్ని} mean? (\textquotedblleft{}Wool\textquotedblright{}.) Say it once more.
\end{tcolorbox}
